% =============================================================================
%  ieee/main.tex -- MINIMAL IEEE-style stub (IEEEtran, two columns)
%
%  This is a stub, not the official template. It reproduces the *shape* of an
%  IEEE conference paper (two columns, IEEE author block, IEEE section order)
%  using only the IEEEtran class, which ships with TeX Live's publishers
%  package. It deliberately does NOT copy any file from IEEE.
%
%  BEFORE SUBMISSION, and before you spend an hour on formatting:
%    1. download the official conference template from
%       https://www.ieee.org/conferences/publishing/templates.html   (verify)
%    2. read the licence/copyright notice that comes with it and keep it in
%       this directory -- do not redistribute the official files from this repo
%    3. copy your text into the official source; layout decisions (page limits,
%       copyright block, keywords) are venue-owned, not yours.
%
%  Compile from THIS directory (relative \input paths depend on it):
%      cd 05_paper/template/ieee && latexmk -pdf main.tex
%  Needs: IEEEtran.cls (Debian package texlive-publishers). If it is missing,
%  build.sh says so instead of failing with 40 lines of TeX errors.
% =============================================================================

\documentclass[conference]{IEEEtran}

\usepackage[utf8]{inputenc}
\usepackage[T1]{fontenc}
\usepackage{amsmath,amssymb}
\usepackage{graphicx}
\usepackage{booktabs}
\usepackage{cite}
\usepackage{url}
\usepackage[hidelinks]{hyperref}

% Two-column layout punishes wide floats; keep tables narrow or use table*.
\newcommand{\method}{TODO-Method-Name}
\newcommand{\fillme}[1]{%
  \par\noindent\fbox{\parbox{0.94\linewidth}{\scriptsize\textbf{TODO} --- #1}}\par}

\begin{document}

\title{TODO(title): specific and falsifiable}

\author{%
\IEEEauthorblockN{TODO(Author One)\IEEEauthorrefmark{1},
                  TODO(Author Two)\IEEEauthorrefmark{1}}
\IEEEauthorblockA{\IEEEauthorrefmark{1}TODO(Department), TODO(Institution)\\
TODO(City), TODO(Country)\\
TODO(email)}
\thanks{%
  \textbf{Use of AI tools.} TODO(disclosure): name each tool and version, the
  sections it touched, and what it did (e.g.\ language editing of
  Section~II). IEEE's policy asks for AI-generated content to be identified by
  section in the acknowledgements; TODO(venue): quote the policy URL and the
  retrieval date. No AI tool produced data, results, figures or references;
  every citation is recorded in \texttt{02\_literature/verified\_citations.csv}.
  The tools are not authors. See Section~\ref{sec:ethics}.}%
}

\maketitle

\begin{abstract}
% TODO(abstract): 150-250 words: problem, method, headline number
% (% num:<id> marker), limitation.
\fillme{Abstract: problem / method / headline result with its index id /
limitation.}
\end{abstract}

\begin{IEEEkeywords}
TODO(keyword 1), TODO(keyword 2), TODO(keyword 3)  % 3-5, lowercase, from the venue's list
\end{IEEEkeywords}

\section{Introduction}
\label{sec:intro}
% TODO(intro): problem / gap / contributions / headline result. IEEE readers
% expect the contribution list as an explicit itemised block.
\fillme{Introduction: problem, gap, contributions, headline result.}

\section{Related Work}
\label{sec:related}
% TODO(related): thematic clusters, then one paragraph on how this differs.
% Only \cite{} keys that exist in 02_literature/verified_citations.csv.
\fillme{Related Work: clusters + delta.}

\section{Method}
\label{sec:method}
% TODO(method): notation, method, cost. One figure at most in two columns.
\fillme{Method: notation, method, cost.}

\section{Experiments}
\label{sec:exp}
% TODO(exp): data, baselines, protocol (seeds, runs, command), metrics.
\fillme{Experiments: data, baselines, protocol, metrics.}

\section{Results}
\label{sec:results}
% TODO(results): answer first, then Table~\ref{tab:main-results}. Numbers are
% generated -- do not retype them.
\fillme{Results: the answer, then the evidence.}

\begin{table}[t]
  \centering
  \caption{TODO(caption): what is compared, on what data, over how many seeds.}
  \label{tab:main-results}
  \IfFileExists{../tables/generated/main-results.tex}{%
    \input{../tables/generated/main-results.tex}%
  }{%
    \fbox{\parbox{0.85\linewidth}{\scriptsize No generated table yet: run
      \texttt{tables/make\_tables.py} (see BUILD.md).}}%
  }
\end{table}

\begin{figure}[t]
  \centering
  \IfFileExists{../figures/generated/main-accuracy.pdf}{%
    \includegraphics[width=\linewidth]{../figures/generated/main-accuracy.pdf}%
  }{%
    \fbox{\parbox[c][2.6cm][c]{0.85\linewidth}{\centering\scriptsize Figure PDF
      not built (needs matplotlib): run \texttt{figures/make\_figures.py}.}}%
  }
  \caption{TODO(caption): axes, unit, uncertainty, and the one thing to notice.}
  \label{fig:main-accuracy}
\end{figure}

\section{Limitations}
\label{sec:limitations}
% TODO(limitations): data, compute, statistics, scope. IEEE does not require
% the section; keep it anyway -- it is the cheapest credibility you can buy.
\fillme{Limitations: data, compute, statistics, scope.}

\section{Ethics Statement}
\label{sec:ethics}
% TODO(ethics): data licences, human subjects (or ``none''), misuse, and the
% AI-assistance statement -- identical text to the title-page footnote.
\fillme{Ethics: licences, misuse, AI-assistance statement.}

\section{Reproducibility Statement}
\label{sec:repro}
% TODO(repro): repo, command, seeds, released artefacts, rerun tolerance.
\fillme{Reproducibility: repo, command, seeds, tolerance.}

\section{Conclusion}
\label{sec:conclusion}
% TODO(conclusion): what is now known, calibrated to the evidence.
\fillme{Conclusion: what is now known, and the next experiment.}

\section*{Acknowledgment}
% TODO(ack): people and funders who helped but are not authors. The AI
% disclosure may live here instead of the footnote if the venue prefers it --
% keep exactly one of the two, never both with different wording.
\fillme{Acknowledgment: non-author help and funding; AI disclosure if the venue
wants it here.}

\bibliographystyle{IEEEtran}
\bibliography{../refs}

\end{document}
